\section{How are probes trained and scored?}
\label{app:probes}

Every probe is a regularised logistic regression on one layer's activation, with the layer chosen inside each training split.
Scores are compared only within a scenario and round, and intervals resample whole groups of related episodes.

\paragraph{Probe recipe.}
The probe is a logistic regression with L2 regularisation ($C = 0.1$), on activations standardised with training statistics only.
Call weights make leaking and non-leaking calls contribute equally within each round, and rounds with only one class get no weight.
We choose the layer from evenly spaced candidates (17 on each model) on an inner validation split of the training data, by the same within-group AUROC that we report.
In the held-out-scenario evaluation, the inner validation set is one of the five training scenarios, so layer selection never sees the held-out scenario.
We run three split seeds and average their AUROCs.

\paragraph{Reasoning probe.}
The reasoning probe uses the same recipe, with every reasoning token carrying its call's label.
Each call's weight is split equally across its tokens, so long reasoning does not dominate training.
A call's score is its maximum token score; the mean is reported as a check (\cref{app:robustness}).

\paragraph{Incoming-only probe.}
The incoming-only probe reads the tokens that are new since the target's previous call, at layer 18 of gpt-oss-20b.
In training, its per-token scores are averaged over sliding 16-token windows before the loss, following \citet{cunningham2026}.
In scoring, each token's score is replaced by a running average that gives weight $2/17$ to the newest token and lets older tokens fade, which spans about 16 tokens.
A call's score is the maximum of this running average over its incoming tokens.

\paragraph{The metric, worked through.}
Suppose one scenario-and-round group holds a leaking call scored 0.8 and two other calls scored 0.5 and 0.9.
The leaking call outranks one of the two, so the group's AUROC is 0.5.
A second group holds a leaking call at 0.7 and one other call at 0.4, so its AUROC is 1.
Each group has one positive, so the reported value is their mean, 0.75.
Groups with only one class are dropped, and ties count one half, which is why any score that is constant within a group gets exactly 0.5.
\Cref{tab:group-coverage} shows, for each headline test, how many leaking calls remain in groups that hold both kinds of call.

\begin{table}[h]
  \centering
  \small
  \setlength{\tabcolsep}{3pt}
  \input{tables/group_coverage.tex}
  \caption{Scenario-and-round groups kept by the metric, and the leaking calls they contain, for each headline test.}
  \label{tab:group-coverage}
\end{table}

\paragraph{Why compare only within scenario and round.}
Leaks fall in earlier rounds than other calls, so rating early rounds as risky already scores well on the raw AUROC.
Within a scenario and round, the same rule scores exactly 0.5.
Leak rates also differ across scenarios (\cref{tab:behaviour-scenario}).

\paragraph{Related episodes and intervals.}
Some episodes share a history, such as variants of one seeded episode and reruns that branch from the same state.
We call each such group a family and always keep a family on one side of a split.
Intervals are 95\% bootstrap intervals: we resample families with replacement 1,000 times, recompute the AUROC from the saved scores without refitting, and take the 2.5th and 97.5th percentiles.

\paragraph{Held-out scenarios.}
We train on five scenarios, test on the sixth, rotate through all six and pool the test scores.
Training also drops any family that appears in the held-out scenario.
Text monitors have no fitted parameters, so their held-out score is their score on the same calls.
\Cref{app:robustness} reports a second evaluation, cross-validation that mixes scenarios.
